\section{Composite Transformations}
\label{sec:composite}

The single-pathway analyses of Sections~\ref{sec:symbolic_internal}--\ref{sec:parametric_centered} treat each carrier transformation as a self-contained unit. A substantial body of recent work, however, chains several such transformations together within a single end-to-end procedure, and in those works the analytical content is rarely located in any individual constituent pathway: it sits in the \emph{integration mechanism} that decides how one pathway's output is converted into the next pathway's input. This section accordingly treats multi-pathway compositions as first-class objects rather than dismantling them across the chapters that handle their individual stages. Four composition patterns are analyzed in turn (\S\ref{sec:composite:extract_then_train}--\S\ref{sec:composite:compress_then_specialize}): three meet a three-paper threshold, while Compress-then-Specialize is retained because its latent intermediate creates a distinct optimization structure. Five further structures with too few representative works to support an independent treatment are registered as candidate patterns (\S\ref{sec:composite:candidate_pool}) and collected together with the four named patterns into a catalog (\S\ref{sec:composite:catalog}). The patterns are built directly on the carrier axis introduced in \S\ref{sec:taxonomy}, with single pathways serving as constructional units rather than as a parallel taxonomic dimension. Each pattern is examined through the same four dimensions used for single pathways---mechanism, cost--benefit structure, failure pathology, and multimodal realization---augmented by its integration mechanism, whose details are exactly what differentiate the four patterns from one another; cross-pattern comparison is reserved for \S\ref{sec:cross_pathway}.

\subsection{Extract-then-Train}
\label{sec:composite:extract_then_train}

The \emph{Extract-then-Train} pattern names the chain Raw $\to$ Refined Tokenized $\to$ Parametric: an LLM-extract operation drawn from the tokenized-internal pathways of \S\ref{sec:symbolic_internal} produces a refined tokenized intermediate, and that intermediate then supervises a downstream policy training run drawn from the tokenized-to-parametric pathways of \S\ref{sec:symbolic_to_parametric}. EvolveR exhibits the canonical two-stage shape: an Offline Self-Distillation phase synthesizes raw trajectories into a repository of reusable strategic principles, and an Online Interaction phase consumes those principles for iterative policy reinforcement. Together with SkillRL, Agentic Proposing, Skill-SD, SaMuLe, Learning from Trials and Errors, and the Skill-driven branch of MetaClaw, EvolveR sits inside a representative set that exceeds the three-paper threshold. Across all of them, the tokenized intermediate functions as a sample-efficiency amplifier---a compact, quality-controlled corpus that replaces a much larger, noisier one---and EvolveR's principle repository, SkillRL's hierarchical SkillBank, and Agentic Proposing's compositional skills are three concrete realizations of that role.

The integration mechanism that holds the two stages together is best read as a coupling between two operations performed on the raw trajectory pool: a filtering operation that decides which trajectories enter the intermediate, and an abstraction operation that compresses the surviving trajectories into higher-order representations. EvolveR performs both---synthesis is the abstraction step and the structure of the repository is the filter---while SkillRL's hierarchy prunes low-quality candidates at each level and SaMuLe's Multi-Level Reflection Synthesis unifies filtering and abstraction across three scales (single-trajectory, intra-task, and inter-task). Skill-library works such as SkillRL and Skill-SD lean toward the abstraction side of this coupling; reflection-driven works such as SaMuLe and Learning from Trials and Errors lean toward filtering. The structural reason the intermediate is doing something direct policy training cannot do becomes visible at this level: gradient-based training weights every example equally and cannot itself discriminate good demonstrations from bad ones in a noisy corpus, so the tokenized intermediate supplies pre-training quality control whose effect is to filter noise before any weight update. EvolveR's structured repository, Agentic Proposing's MGPO, and Skill-SD's importance-weighted reverse-KL loss are three variants of this same pre-training-quality-control machinery; the intermediate is structurally necessary when raw-trajectory noise is high, as it typically is for self-generated rollouts on long-horizon tasks, and is correspondingly less central for clean expert traces where direct training already works well.

The cost--benefit structure follows directly: Extract-then-Train spends extra extraction compute in exchange for a denoised training corpus, and the trade turns on the noise floor of the raw experience pool. The pattern's distinctive failure mode is intermediate-quality propagation---errors introduced during extraction become ground-truth supervision for the downstream training run and are therefore amplified rather than absorbed---and this pathology is structurally absent from either constituent pathway when run alone. Multimodal instances concentrate in GUI and embodied settings, with Learning from Trials and Errors' dual reflection-in-action and reflection-on-action mode the most representative case. SkillRL's recursive evolution and EvolveR's reported gains both suggest that the composite delivers value beyond what its constituent stages would deliver in isolation, but no controlled ablation in the current literature isolates the intermediate's contribution from the rest of the pipeline, so evidence for non-additive gains remains suggestive rather than conclusive.

\subsection{Evaluate-then-Optimize}
\label{sec:composite:evaluate_then_optimize}

The \emph{Evaluate-then-Optimize} pattern names the chain Raw $\to$ Evaluator $\to$ Policy: a parametric evaluator trained on tokenized experience supplies the gradient signal that drives policy optimization. As already observed in \S\ref{sec:p6}, this pattern is where P6 actually appears in the literature; standalone P6 is structurally rare, and what differentiates the works in this group is precisely that evaluator construction and policy update are presented as a joint contribution rather than as two independent ones. SWEET-RL, From Novice to Expert, ReST-MCTS, and PRL together form the representative set. The evaluator's role here shifts in kind from its standalone P4 use: scores no longer serve only to rank candidate outputs but become advantages that become gradient signals, a use absent from inference-only Score-and-Select. ReST-MCTS illustrates the dual function most cleanly, with the same inferred rewards serving simultaneously as PRM training targets and as filters that select which traces enter the policy training set.

The integration mechanism in this pattern is the conversion path from evaluator output to a gradient signal that policy training can consume, and the four representative works diverge precisely on what gets backpropagated and where. SWEET-RL converts critic step-rewards into DPO preference pairs; From Novice to Expert folds discriminator output into PPO per-step rewards via implicit-reward and inverse RL; ReST-MCTS uses one PRM output for two purposes, training the PRM itself against value targets and selecting high-quality traces for downstream policy training; PRL decomposes outcome reward into per-step supervision through entropy-regularized RL. The four mechanisms differ in routing but share the structural choice of treating the evaluator as a gradient-shaping mechanism rather than as a passive scorer. The reason the evaluator is doing work that outcome-based RL cannot do becomes visible here: a single reward per trajectory makes per-step credit assignment intractable for long horizons, and the evaluator supplies the per-step credit that outcome decomposition cannot generate. ReST-MCTS's tree-search reward inference, which estimates per-step values from rollout success rates, produces exactly the kind of dense signal that outcome reward alone cannot deliver, and the composite is correspondingly justified in regimes where credit assignment dominates evaluator training cost---multi-turn reasoning in SWEET-RL, multi-step mathematics in ReST-MCTS, multi-step inference in PRL---but is unnecessary on short-horizon tasks where outcome supervision suffices.

Read as a unit, Evaluate-then-Optimize adds the evaluator-training cost analyzed in \S\ref{sec:p4} in exchange for step-level credit signals that would otherwise be unavailable; the implicit-PRM line discussed in \S\ref{sec:p4_process} lowers this cost and shifts the break-even point. The pattern's distinctive failure is evaluator-bias propagation: bias enters the policy weights through the gradient and becomes permanent, in a way that standalone P4 use does not produce, where bias affects only inference-time selection. ReST-MCTS's concurrent PRM refinement and SWEET-RL's asymmetric-information critic are partial mitigations rather than structural cures. Multimodal instances are sparse---PROPA's visual RL is the principal case---and the sparsity reflects the upstream scarcity of multimodal P4 evaluators rather than any pattern-specific obstacle. Reported gains over the natural single-pathway baselines are large enough to suggest non-additive value, with ReST-MCTS's dual-purpose design coming closest to a structural argument, but as in Extract-then-Train no work in the representative set isolates the evaluator's contribution from the specific architectural choices that surround it.

\subsection{Self-Generation Loop}
\label{sec:composite:self_generation}

The \emph{Self-Generation Loop} pattern names the closed-loop chain Policy $\to$ Raw $\to$ Policy: a policy generates trajectories, a filter selects a subset, and that subset is used to train the same policy, with the cycle iterated. Beyond Human Data and ReST$^{\text{EM}}$ established the canonical form of the loop---generate, filter with binary feedback, fine-tune, and repeat---while Self-Instruct introduced the data-side variant. The representative set here is the largest in this section, exceeding nineteen works and including Agent-R, Agent0, Beyond Policy Optimization, Internalizing Agency, START, EvoCUA, MobileGUI-RL, AgentEvolver, and the system-level integrations within UI-TARS-2 and V-Zero. The defining structural feature is that the policy occupies two roles simultaneously, as data generator and as data consumer, in a way no single pathway in isolation can produce. The contrast with a one-shot composition of P7 and P5 is informative: that one-shot chain produces a static dataset, while the loop produces an evolving dataset that adapts to the moving policy distribution and is the proper unit of analysis only when treated as a unit.

The integration mechanism in this pattern is the filter that decides which generated trajectories enter the next iteration, and filter design partitions the literature into four recognizable types whose choices determine collapse risk, diversity preservation, and per-iteration compute cost. The simplest type is outcome-based, exemplified by Beyond Human Data, which retains trajectories on binary success and is correspondingly the cheapest and the weakest against collapse. The process-based type, exemplified by Agent-R, uses MCTS to recover correct trajectories from erroneous ones and pairs the search with a model-guided critique that identifies the first error step, so the filter consumes step-level diagnostics rather than only outcome signals. The LLM-based type, exemplified by Internalizing Agency, lets self-reflection on environmental feedback decide what to keep or revise. The asymmetric self-play type, exemplified by Agent0, places a curriculum agent that proposes frontier tasks against an executor agent that learns to solve them, so the filter pressure comes from an adversary rather than from a static criterion. The reason the loop is doing work a single-shot P5 with a static dataset cannot do is visible here: as policy training progresses, the optimal training distribution shifts, and data that is informative at iteration $k$ may be uninformative at iteration $k{+}10$; a static dataset cannot track this shift, while regenerated data can. The composite is correspondingly justified when distribution shift during training is non-negligible---the typical condition for long-horizon, exploratory, and open-ended tasks---and unnecessary in regimes where the training distribution is stable.

Read as a unit, the loop multiplies compute by iteration count but eliminates per-iteration annotation cost, and becomes economically attractive when high-quality external data is expensive. Three structural failures are inherent to the closed-loop topology rather than to any specific filter choice. Model collapse describes the gradual degradation of output quality when self-training feeds on its own generations across many iterations. Distribution collapse describes the narrowing of the policy's output distribution as the filter selects trajectories increasingly similar to one another. Diversity loss describes the bound that the loop's reachable distribution is fixed by the initial policy and admits no external novelty. Better filter design delays each of these but cannot escape any of them, because the loop's sole information source is the policy itself. Multimodal loops concentrate in GUI agents---UI-TARS-2's data flywheel, MobileGUI-RL's trajectory-aware advantages, EvoCUA's verifiable synthesis---with the additional requirement that visual state realism be verifiable. Agent0's reported gains in the absence of external data and Beyond Human Data's favorable scaling jointly suggest that iterative self-training extracts capability beyond what linear data scaling would predict, but as elsewhere in this section no published ablation isolates the loop's contribution from the surrounding architectural choices, so evidence for non-additive gains remains suggestive rather than conclusive.

\subsection{Compress-then-Specialize}
\label{sec:composite:compress_then_specialize}

The \emph{Compress-then-Specialize} pattern names the chain Raw $\to$ Latent $\to$ Parametric: raw experience is compressed into a latent representation that subsequently drives downstream parametric training. LatentEvolve is the only standalone instance in the current literature: its daytime scaling phase retrieves historical latents to guide reasoning, and its nighttime scaling phase integrates accumulated latent optimizations into a parametric refiner. Although this pattern does not meet the three-paper threshold used for the other named patterns, it is retained because a latent intermediate enables a form of continuous-space joint optimization that a tokenized intermediate cannot in principle support. Viewed along the same dimensions as single pathways, the pattern's distinctive property is that the latent intermediate is continuously differentiable, so end-to-end backpropagation can flow through it in a way that any tokenized intermediate would necessarily break. This trades editability against end-to-end trainability relative to Extract-then-Train, and its principal risk---latent degeneration, in which the intermediate collapses into a thin parametric abstraction---is what LatentEvolve's daytime--nighttime alternation is designed to forestall. Multimodal instances are currently absent from the literature.

The integration mechanism is bidirectional compression-and-retention: daytime latents guide reasoning, the resulting reasoning outcomes update the latent store, and that store in turn updates the parametric refiner. The reason for retaining this pattern despite its single-instance literature follows from the fact that continuous representations retain gradient connectivity across the full pipeline and therefore admit end-to-end training against a downstream objective; a tokenized intermediate severs this gradient, so the extraction operation itself cannot be optimized for downstream performance. This differentiability advantage is the central argument for treating Compress-then-Specialize as a distinct pattern, and it is independent of instance count. The single-instance literature does not support an empirical claim of non-additive performance gain; the distinct property that does apply---end-to-end differentiability unavailable to Extract-then-Train---is a design property rather than a measured effect, and is reported as such.

\subsection{Candidate Patterns}
\label{sec:composite:candidate_pool}

Five further composition structures fall short of the three-paper threshold under the current literature but are recognizable enough to be registered with their link shape, current representatives, and the specific gap between current coverage and threshold. \emph{Formalize-then-Optimize} (N-Tok $\to$ S-Tok $\to$ $\pi$-Par) is currently represented by ScoreFlow and Meta Context Engineering, leaving a gap of one. \emph{Externalize-for-Training} ($\pi$-Par $\to$ N-Tok $\to$ V-Par), in which a policy externalizes data to train an evaluator, is represented by Self-Improving VLM Judges and is structurally distinct from both Self-Generation Loop and standalone P7; the gap is two. \emph{Internalize-then-Externalize} (N-Tok $\to$ $\pi$-Par $\to$ N-Tok skill), whose topology is the reverse of Extract-then-Train, is represented by SAGE and RL for Self-Improving Agent with Skill Library and is the most likely candidate for promotion, with a gap of one. \emph{Co-Evolving World Model} ($\pi$-Par $\leftrightarrow$ V-Par), in which the world model is an independent parametric component structurally distinct from policy-as-generator and evaluator-as-signal, is represented by WebEvolver, WebSynthesis, and a partial contribution from SEAgent, making it near-threshold but not yet a fully supported pattern. \emph{Neural-Symbolic Self-Learning}, exemplified by SymAgent, sits at the boundary of standard agent literature and falls short by two. All five candidates are revisited in \S\ref{sec:open_problems}.

\subsection{Pattern Catalog}
\label{sec:composite:catalog}

The four analyzed patterns and the five candidate entries together constitute the composition catalog used as the navigation anchor for \S\ref{sec:cross_pathway}; cross-pattern trade-off comparisons are deferred to \S\ref{sec:cross_pathway:trade_off}, and the structural rules that govern admissible compositions are deferred to \S\ref{sec:cross_pathway:grammar}.
